\documentclass[a4paper]{article}
\usepackage{amsfonts}
\usepackage{amsmath}
\usepackage{amssymb}
\usepackage{graphicx}

\begin{document}
  
\noindent \large Auteur: Michael Livchits \\
\noindent \large Studierichting: Informatica\\
\noindent \large Oefeningenreeks 6B nr. 21.3\\

\medskip

\normalsize

$\dfrac{2(2n+3)}{n(n+1)(n+2)(n+3)} = \dfrac{A}{n} +\dfrac{B}{n+1}+\dfrac{C}{n+2}+\dfrac{D}{n+3}$\\

$A = \dfrac{6}{6} = 1$\\

$B = \dfrac{2}{-2} = -1$\\

$C = \dfrac{-2}{2} = -1$\\

$D = \dfrac{-6}{-6} = 1$\\

$\dfrac{2(2n+3)}{n(n+1)(n+2)(n+3)} = \dfrac{1}{n} -\dfrac{1}{n+1}-\dfrac{1}{n+2}+\dfrac{1}{n+3}$\\

Eerste 4 termen opschrijven:\\

$s_1 = \dfrac{1}{1}-\dfrac{1}{2}-\dfrac{1}{3}+\dfrac{1}{4}$\\

$s_2 = \dfrac{1}{2}-\dfrac{1}{3}-\dfrac{1}{4}+\dfrac{1}{5}$\\

$s_3 = \dfrac{1}{3}-\dfrac{1}{4}-\dfrac{1}{5}+\dfrac{1}{6}$\\

$s_4 = \dfrac{1}{4}-\dfrac{1}{5}-\dfrac{1}{6}+\dfrac{1}{7}$\\

$s_{n-1} = \dfrac{1}{n-1}-\dfrac{1}{n}-\dfrac{1}{n+1}+\dfrac{1}{n+2}$\\

$s_n = \dfrac{1}{n}-\dfrac{1}{n+1}-\dfrac{1}{n+2}+\dfrac{1}{n+3}$\\

Bij optellen zien wij duidelijk dat bijna alle termen afvallen behalve:\\

$\Rightarrow 1-\dfrac{1}{3}-\dfrac{1}{n+1}+\dfrac{1}{n+3}$\\

$= \dfrac{3(n+3)(n+1)-(n+3)(n+1)-3(n+3)+3(n+1)}{3(n+3)(n+1)}$\\

$= \dfrac{2(n+3)(n+1)-3(n+3)+3(n+1)}{3(n+3)(n+1)}$\\

$= \dfrac{2n^2+2n+6n+6-3n-9+3n+3}{3(n+3)(n+1)}$\\

$= \dfrac{2n^2+8n}{3(n+3)(n+1)}$\\

\begin{thebibliography}{999}
%\bibitem{a} Referentie
\end{thebibliography}
\end{document} 